\section{An independent coupling proof of the square-root bound}
\label{proof:app:coupling}

This appendix proves Proposition~\ref{proof:prop:squareroot} for the
cut chain without the harmonic decomposition. It provides a
probabilistic verification of the square-root estimate. The constants
in the resulting total-variation bounds may depend on $n$ and
$\beta$; only their exponential rate is used.

Assume $n\geq8$, $4/n^2\leq\beta\leq1/16$, and put
$a=n\sqrt\beta\in[2,n/4]$. Start two copies of the cut process from
adjacent nonempty sets. Until they meet, the coupling has two stages.
After meeting, run the copies identically.

\subsection*{Stage I: one set contains the other}
Suppose $X\subset Y$, $Y\setminus X$ consists of one bond, and
$|X|=k$. Add every common vacant bond to both sets at rate $\beta$.
For each common occupied bond, remove it from both sets at rate
$(k-1)/n$, and remove it from $Y$ alone at rate $1/n$. At the unique
bond in $Y\setminus X$, add it to $X$ at rate $\beta$ and remove it
from $Y$ at rate $k/n$.

The first marginal has birth rate $\beta$ per vacant bond and death
rate $(k-1)/n$ per occupied bond. The second has birth rate $\beta$
per vacant bond and death rate $k/n$ per occupied bond. Hence both
marginals agree with~\eqref{proof:eq:LB}. The last two events couple
the copies; a removal from $Y$ alone at a common bond starts
Stage II. Until one of these exits, $k$ increases at rate
$\beta(n-k-1)$ and decreases at rate $k(k-1)/n$. The total exit rate
is $\beta+2k/n$. Thus the killed internal generator is exactly
$\mathcal A$ in~\eqref{proof:eq:derivativegenerator}.

\subsection*{Stage II: equal sizes and one exclusive bond each}
Write $|X|=|Y|=r+1$, so the intersection contains $r$ bonds. Add
each common vacant bond to both sets at rate $\beta$, and remove
each common occupied bond from both at rate $r/n$. At rate
$\beta$, simultaneously add to each set the bond exclusive to the
other set, so both become their union. At rate $r/n$,
simultaneously remove each set's exclusive bond, so both become
their intersection. When $r=0$, the latter event has rate zero and
does not create the empty set.

Each simultaneous absorption event produces one jump in each
marginal. Accordingly the total absorption rate is $\beta+r/n$.
The internal coordinate $r\in\{0,\ldots,n-2\}$ has killed generator
\begin{align}
 (\mathcal Bh)(r)
 ={}&\beta(n-r-2)(h(r+1)-h(r))
     +\frac{r^2}{n}(h(r-1)-h(r))\notag\\
 &-\left(\beta+\frac rn\right)h(r).
 \label{proof:eq:stageII}
\end{align}
The marginal birth and death rates are again those
of~\eqref{proof:eq:LB}.

\subsection*{Positive certificates and coupling time}
For both stages use $h(z)=1/(z+a)$. For Stage I,
\eqref{proof:eq:countreciprocal} gives
$-(\mathcal Ah)(k)/h(k)\geq\sqrt\beta/5$. For Stage II, exact
cancellation yields
\begin{equation}\label{proof:eq:stageIIcertificate}
 \frac{-(\mathcal Bh)(r)}{h(r)}
 =\frac{\beta(n-r-2)}{r+a+1}+\beta
  +\frac{r(a-1)}{n(r+a-1)}.
\end{equation}
If $r\leq a$, then $n-r-2\geq n/2$ and $r+a+1\leq5a/2$, so
the first term is at least $\sqrt\beta/5$. If $r\geq a$, then
$r/(r+a-1)\geq1/2$ and $a-1\geq a/2$, so the last term is at
least $a/(4n)=\sqrt\beta/4$. Consequently, with
$c=\sqrt\beta/5$, both killed generators satisfy $Ah\leq-ch$.

For completeness, any finite matrix $A$ with nonnegative
off-diagonal entries generates a positivity-preserving semigroup:
add a sufficiently large scalar multiple of the identity to make
the whole matrix nonnegative and use the exponential series. If
$h>0$ and $Ah\leq-ch$, differentiation gives
\[
 \frac{d}{dt}\bigl(e^{t(A+cI)}h\bigr)
 =e^{t(A+cI)}(A+cI)h\leq0,
\]
and hence $e^{tA}h\leq e^{-ct}h$. The survival time $T_i$ of each
stage therefore satisfies, uniformly over its starting states,
\[
 \mathbb P(T_i>t)\leq C_i e^{-ct},
 \qquad
 C_i=\frac{\max h}{\min h}<\infty.
\]
For every $0<\theta<c$, integration of this tail bound shows that
\[
 \mathbb E e^{\theta T_i}
 \leq1+\frac{\theta C_i}{c-\theta}.
\]
The strong Markov property at entry to Stage II, using its uniform
bound over initial states, gives a finite exponential moment at
rate $\theta$ for the total coupling time. Stage II has duration
zero if Stage I couples directly. Thus the total coupling time has
tail $O(e^{-\theta t})$.

The coupling inequality gives this rate for total variation between
chains started at any adjacent nonempty sets. Any two nonempty
sets are connected by a finite path through their union, using
single-bond additions and deletions and never visiting the empty
set. The triangle inequality gives the same rate for any two
initial states.

Finally, let $f$ be a nonconstant eigenfunction of $L_B$ with
eigenvalue $-\lambda$, and choose $x,y$ with $f(x)\neq f(y)$. Then
\[
 e^{-\lambda t}|f(x)-f(y)|
 =|\mathbb E_x f(B_t)-\mathbb E_y f(B_t)|
 \leq2\|f\|_\infty\,
       \|P_t(x,\cdot)-P_t(y,\cdot)\|_{\mathrm{TV}}.
\]
The established exponential bound forces $\lambda\geq\theta$.
Since this holds for every $\theta<c$, it follows that
$\gamma_B\geq\sqrt\beta/5$. Proposition~\ref{proof:prop:comparison}
then gives the polymer bound.
